\section{Authenticated Message-Oracle Implementations}

The algebraic protocols below need authenticated access to fixed tensors.  The
ideal interface is not a second ideal service; it is the message
oracle from Figure~\ref{fig:func-msg-oracle}.  A prover, an indexer, a Merkle
tree, or a virtual construction fixes a logical tensor and later answers
\(\mathcal F_{\mathrm{MsgOracle}}.\mathsf{Open}\) queries at named addresses.

The committed object carries its own natural address space.  If that address
space is \(\B^d\), the Merkle implementation uses the natural binary leaf order
\[
  \bits_d(0),\bits_d(1),\ldots,\bits_d(2^d-1).
\]
If the address space is a product of binary domains, the leaf order is the
natural binary order after concatenating the address blocks.  Thus the
message-oracle API does not hide any indexing convention: the logical address
space is the layout.  In the codeword handles below, the value space is \(\F\).
If \(\mathcal F_{\mathrm{MsgOracle}}.\mathsf{Open}(R,a)\ne\bot\), its value is
the unique logical value bound to \(R\) at address \(a\).  Equivalently, a
verifier can check a claimed value \(v\) by opening \(R\) at \(a\) and checking
that the returned value equals \(v\).

The ordinary implementation is a Merkle tree.  The box below is not a new ideal
functionality; it is one concrete way to implement the
\(\mathcal F_{\mathrm{MsgOracle}}.\mathsf{Open}\) interface.  For simplicity it
is written for address space \(\B^d\).  Product address spaces are first
serialized in the natural binary order described above.

\begin{protocolbox}{Implementation: Merkle Message Oracle}{fig:impl-merkle-msg-oracle}
\noindent Public parameters: a collision-resistant hash function
\(\mathsf{Hash}\) and a value encoder into hash inputs.

\smallskip
\noindent\(\underline{\Pi_{\mathrm{MerkleMsgOracle}}.\mathsf{Commit}(\prover:x\in\mathcal L^{\B^d},\ \verifier:\bot)\to R}\)

\smallskip
\begin{enumerate}[leftmargin=0.24in,label=\textbf{M\arabic*.},itemsep=0.08em,topsep=0.1em,parsep=0pt]
  \item \(\prover\): for \(i\in\{0,\ldots,2^d-1\}\), set
        \(T_0[i]\gets\mathsf{Hash}(x[\bits_d(i)])\).
  \item \(\prover\): for \(\ell\in\{0,\ldots,d-1\}\) and
        \(j\in\{0,\ldots,2^{d-\ell-1}-1\}\), set
        \[
          T_{\ell+1}[j]
          \gets
          \mathsf{Hash}(T_\ell[2j]\Vert T_\ell[2j+1]).
        \]
  \item \(\prover\): set \(R\gets T_d[0]\).
  \item \(\prover\to\verifier\): \(R\).
  \item \(\prover,\verifier\): return \(R\).
\end{enumerate}

\smallskip
\noindent\(\underline{\Pi_{\mathrm{MerkleMsgOracle}}.\mathsf{Open}(\verifier:(R,a=\bits_d(i)),\ \prover:(x,a))\to(\verifier:o\in\mathcal L\cup\{\bot\})}\)

\smallskip
\begin{enumerate}[leftmargin=0.24in,label=\textbf{O\arabic*.},itemsep=0.08em,topsep=0.1em,parsep=0pt]
  \item \(\prover\): send \(v\gets x[a]\).
  \item \(\prover\): set \(j_0\gets i\).  For each
        \(\ell\in\{0,\ldots,d-1\}\), send the sibling
        \(s_\ell\gets T_\ell[j_\ell\oplus1]\) and set
        \(j_{\ell+1}\gets\lfloor j_\ell/2\rfloor\).
  \item \(\verifier\): set \(h_0\gets\mathsf{Hash}(v)\) and \(j_0\gets i\).
  \item \(\verifier\): for each \(\ell\in\{0,\ldots,d-1\}\), compute
        \[
          h_{\ell+1}\gets
          \begin{cases}
            \mathsf{Hash}(h_\ell\Vert s_\ell), & j_\ell\equiv0\pmod 2,\\
            \mathsf{Hash}(s_\ell\Vert h_\ell), & j_\ell\equiv1\pmod 2,
          \end{cases}
        \]
        and set \(j_{\ell+1}\gets\lfloor j_\ell/2\rfloor\).
  \item \(\verifier\): if \(h_d=R\), return \(v\); otherwise return \(\bot\).
\end{enumerate}
\end{protocolbox}

We will use several implementations of this same API.  The ordinary
implementation is Merkle: it stores the logical tensor itself and opens
coordinates by Merkle paths.  Blaze also uses virtual handles: a virtual handle
does not store one flat tensor directly, but it still implements
\(\mathcal F_{\mathrm{MsgOracle}}.\mathsf{Open}\).  The names
\(\mathsf{RowFold.Commit}\) and \(\mathsf{CompositeCommit}\) use
\(\mathsf{Commit}\) in this constructor sense: they create a handle.  They do
not introduce a second coordinate-opening API.  After the handle exists, the
only coordinate operation is \(\mathsf{Open}\).

The main virtual implementation used by Blaze is
\(\mathsf{RowFold.Commit}\), which folds the committed row-word matrix into
the folded-tensor handle \(R_\star\).  The later \(\MLEAuth\) implementation
also constructs packed coordinate handles by routing lane openings to ordinary
or virtual message-oracle handles.

For example, if $R_c$ is a handle to a column tensor
$C\in(\F^{\B^\ell})^{\B^n}$ and $\rho\in\F^{\B^\ell}$, then a compressed
\(\mathsf{RowFold.Commit}(R_c,\rho)\) returns a virtual
handle \(R_\star\) implementing
\(\mathcal F_{\mathrm{MsgOracle}}.\mathsf{Open}\).  Its logical domain is \(\B^n\) and
its logical tensor is
\[
  c_\star[b]=\ip{\rho}{C[b]}.
\]
Opening $R_\star$ at $b$ is the existing Open interface: it opens $R_c$ at
$b$, receives the already-bound column value \(w=C[b]\in\F^{\B^\ell}\), and
returns \(\sum_{h\in\B^\ell}\rho[h]w[h]\).  The backend sees only the handle
$R_\star$ and the operation
\(\mathcal F_{\mathrm{MsgOracle}}.\mathsf{Open}(R_\star,b)\).

The packed \(\MLEAuth\) handle is virtual in the same sense.  It is a handle
over a lane domain \(\mathcal L\times\B^D\).  Opening a lane routes to the
coordinate handle supplied for that lane, which may itself be ordinary,
composite, or virtual.  The backend sees only the packed handle, not the
machinery behind each lane.

We will also use composite virtual handles.  This is implementation notation
for \(\mathcal F_{\mathrm{MsgOracle}}.\mathsf{Open}\), not a new ideal
functionality or cryptographic assumption.  A composite handle stores a public
routing table between existing handles over disjoint coordinate sets.  If
\(E_1,\ldots,E_m\) are disjoint address sets and component handles \(R_i\)
open tensors over \(E_i\), then
\[
  R=\mathsf{CompositeCommit}((R_1,E_1),\ldots,(R_m,E_m))
\]
denotes the virtual handle over \(D=E_1\sqcup\cdots\sqcup E_m\).  Opening
\(R\) at \(a\in D\) finds the unique \(i\) with \(a\in E_i\) and returns
\(\mathcal F_{\mathrm{MsgOracle}}.\mathsf{Open}(R_i,a)\).

For the systematic split used by BaseFold, the message addresses are
\(\mathcal L\times\B^k\).  The codeword addresses are
\(\mathcal L\times\B^n\).  The systematic address of
\((q,a)\in\mathcal L\times\B^k\) is \((q,0^{n-k}\Vert a)\).  Thus
\[
  D_{\mathrm{sys}}
  =
  \{(q,0^{n-k}\Vert a):q\in\mathcal L,\ a\in\B^k\},
  \qquad
  D_{\mathrm{ext}}
  =
  \mathcal L\times\bigl(\B^n\setminus\bits_n([0,2^k-1])\bigr).
\]
The composite handle is
\[
  \mathsf{CompositeCommit}
  ((C_A,D_{\mathrm{sys}}),(C_{\mathrm{ext}},D_{\mathrm{ext}})).
\]
Here \(C_A\) is viewed as a handle over \(D_{\mathrm{sys}}\): its value at
\((q,0^{n-k}\Vert a)\) is the message value \(A_q[a]\).  The composite handle
opens systematic codeword addresses through \(C_A\), and opens the remaining
codeword addresses through \(C_{\mathrm{ext}}\).

The handle is binding at the logical-tensor level.  It is not an algebraic
proof: it does not by itself prove that the logical tensor is a codeword or
that it satisfies any relation.  The algebraic proof comes from random queries
and low-degree/folding checks layered on top of these authenticated openings.
